\RequirePackage{silence}
\WarningsOff % Turn off all warnings from all packages
\ErrorsOff   % (Caution) Use only if you are sure the code has no severe logic errors
\documentclass[a4paper,14pt, twoside]{memoir}
\usepackage[utf8x]{inputenc}
\usepackage[english]{babel}    
\usepackage{url}
\usepackage{microtype}
\usepackage{color}
\usepackage{amsmath}
\usepackage{amsthm}         % Provides proof environment
\usepackage{amssymb}
\usepackage{mathtools}     % Includes amsmath but optimized and with more features
\usepackage{hyperref}
\usepackage{enumitem}
\usepackage{autobreak}
\usepackage{array}          % Support for notation table formatting
\usepackage{booktabs}       % Professional table rules

\hypersetup{
	colorlinks=true,
	linkcolor=blue,
	filecolor=magenta,
	urlcolor=cyan
}

% Theorem Environments
\newtheorem{definition}{Definition}[chapter]
\newtheorem{md}{Proposition}[chapter]
\newtheorem{theorem}{Theorem}[chapter]
\newtheorem{proposition}{Proposition}[chapter]
\newtheorem{corollary}{Corollary}[chapter]
\newtheorem{remark}{Remark}[chapter]
\newtheorem{example}{Example}[chapter]

% Declare Math Operators
\DeclareMathOperator{\Tr}{Tr}
\DeclareMathOperator{\Sym}{Sym}
\DeclareMathOperator{\setu}{s}
\DeclareMathOperator{\dimU}{\dim}
\DeclareMathOperator{\w}{w}
\DeclareMathOperator{\card}{card}
\DeclareMathOperator{\pr}{pr}
\DeclareMathOperator{\pgr}{pgr}
\DeclareMathOperator{\scale}{scale}
\DeclareMathOperator{\rad}{rad}
\DeclareMathOperator{\Val}{Val}
\DeclareMathOperator{\Conf}{Conf}

% Custom commands
\newcommand{\set}[1]{\left\{#1\right\}} % set
\newcommand{\tbd}[1]{\mathcal{#1}} % uncertain set
\newcommand{\s}[1]{\set{#1}_\mathrm{u}} % uncertain number
\newcommand{\kbd}[1]{{#1}_\mathrm{u}} % uncertain interval
\newcommand{\bd}[1]{\left(#1\right)_\mathrm{u}} % uncertain expression
\newcommand{\tsbd}[1]{\mathbf{U}(\mathbb{#1})}
\newcommand{\tsbdcapn}[2]{\mathbf{U^{#2}}(\mathbb{#1})}
\newcommand{\tsbdts}[1]{\mathbf{U}_{\w}(\mathbb{#1})} % set of uncertain numbers
\newcommand{\ddclmd}{\mathcal{V}} % truth measure of a proposition
\newcommand{\ts}[1]{\w_{#1}} % Weight
\newcommand{\gtcl}[2]{\left(#1\right)_{#2}} % truth measure
\newcommand{\gr}[1]{\text{gr}(#1)} % graph of a function
\newcommand{\hpw}[2]{#1_1\left(#2\right)} % point-wise function
\newcommand{\hbd}[2]{#1_m\left(#2\right)} % uncertain function
\newcommand{\eq}[1]{\left[#1\right]_{eq}} % eqs operator
\newcommand{\kgdd}[1]{\left(#1\right)_1} % pointwise operator
\newcommand{\kgddmr}[1]{\left(#1\right)_{1'}} % real operator
\newcommand{\kgdt}[1]{\left(#1\right)_m} % pointwise operator
\newcommand{\kgdtmr}[1]{\left(#1\right)_{m'}} % pointwise operator

\title{Uncertain Set Theory}
\author{
	Tran Manh Cuong\thanks{Email: \href{mailto:tmcuong2408@gmail.com}{tmcuong2408@gmail.com} | Phone: 0353237140} \\[0.5em]
	\small Alumnus of Faculty of Mathematics and Informatics (K11)
	\\ Thai Nguyen University of Sciences \\
	\small Address: DJ7 Street, Thoi Hoa, Ho Chi Minh City, Vietnam \\
	\small \textit{Research Field: Mathematics of Uncertainty}
}
\date{Last updated: \today}

\begin{document}
	
	\maketitle
	
	\begin{abstract}
		This paper establishes \textbf{Uncertain Set Theory}, a novel foundational mathematical framework built from First Principles to model uncertainty without relying on probabilistic distributions, fuzzy membership degrees, or static interval bounds. Traditional approaches such as Interval Arithmetic frequently suffer from correlation loss and catastrophic numerical noise explosion, while probabilistic models require prior knowledge of probability densities that are often unavailable in real-world high-noise environments. 
		
		We address these fundamental limitations by introducing the concepts of \textbf{Uncertain Elements} and \textbf{Uncertain Sets}, defined as multi-state superposition entities characterized by possibility domains $\mathrm{Val}(x)$ and configuration spaces $\mathrm{Conf}(S)$. Through this scenario-based structural formulation, we demonstrate that set cardinality itself unifies naturally as an uncertain number $|S| \in \mathbf{U}(\mathbb{N})$, resolving cardinality degeneration and state-collapse seamlessly. Furthermore, algebraic operations on uncertain sets inherit classical binary set relations over their configuration spaces, preserving absolute logical correlation and enabling canonical complexity reduction down to $\mathcal{O}(1)$ via Lazy Evaluation. This framework serves as the core foundation for $U(\mathbb{K})$ arithmetic, scenario-space optimization, and high-dimensional uncertain topology.
		
		\vspace{1em}
		\noindent \textbf{Keywords:} Uncertain Set Theory, Uncertain Elements, Configuration Space, State Superposition, Uncertain Cardinality, Scenario-Based Arithmetic, Correlation Loss, $U(\mathbb{K})$ Spaces, First Principles.
	\end{abstract}
	
	\newpage
	\tableofcontents
	\newpage
	\chapter{Uncertain Sets}
	
	\begin{definition}[Uncertain Element]
		Let $U$ be a non-empty universal set. An \textbf{uncertain element} $x$ on $U$ is an object taking an undetermined value in $U$, constrained by a set of possible values $A \subseteq U$ ($A \neq \emptyset$), denoted by:
		\[
		x \in A
		\]
		The set $A$ is called the \textbf{possibility domain} (or state space) of $x$, and we denote $\mathrm{Val}(x) = A$.
	\end{definition}
	
	\begin{definition}[Uncertain Set]
		Let $U$ be a non-empty universal set and $I$ be an index set. An \textbf{uncertain set} $S$ on $U$ is a collection of uncertain elements of the form:
		\[
		S = \{x_i : i \in I\}
		\]
		where each $x_{i}$ ($i \in I$) is an uncertain element on $U$ with possibility domain $A_i = \mathrm{Val}(x_{i}) \subseteq U$.
		
		A \textbf{realized configuration} (or empirical state) of $S$ is a classical set $S^* \subseteq U$ obtained when the elements $x_{i}$ simultaneously take specific values $v_i \in A_i$:
		\[
		S^* = \{v_i : i \in I\}
		\]
		The collection of all possible configurations of $S$ is denoted by $\mathrm{Conf}(S)$.
	\end{definition}
	
	\begin{definition}[Uncertain Cardinality]
		Let $S = \{x_i : i \in I\}$ be an uncertain set on the universal set $U$. The \textbf{uncertain cardinality} (or number of uncertain elements) of $S$, denoted by $|S|$, is an uncertain number on the set of natural numbers $\mathbb{N}$ whose possibility domain is defined by:
		\[
		\mathrm{Val}(|S|) = \{ |S^*| : S^* \in \mathrm{Conf}(S) \} \subseteq \mathbb{N}
		\]
		where $|S^*|$ represents the classical cardinality of the configuration $S^*$.
	\end{definition}
	
	\begin{remark}
		If all uncertain elements $x_i$ in $S$ take distinct values across their entire possibility domains, we have $|S^*| = |I|$ for all $S^* \in \mathrm{Conf}(S)$, in which case the uncertain cardinality degenerates into a classical number $|S| = |I|$. Conversely, size degeneration (cardinality collapse) occurs at configurations $S^*$ where duplicate values appear among the elements ($v_i = v_j$ for $i \neq j$).
	\end{remark}
	
	\begin{definition}[Binary Operations on Uncertain Sets]
		Let $S_1$ and $S_2$ be two uncertain sets on the universal set $U$, and let $\star$ be a classical binary set operation (such as $\cup, \cap, \setminus$). The corresponding binary operation on uncertain sets, denoted by $S_1 \star S_2$, is an uncertain set on $U$ whose configuration space is defined by:
		\[
		\mathrm{Conf}(S_1 \star S_2) = \{ S_1^* \star S_2^* : S_1^* \in \mathrm{Conf}(S_1), \, S_2^* \in \mathrm{Conf}(S_2) \}
		\]
	\end{definition}
	
	\begin{example}[Uncertain Union and Intersection]
		From the general definition above, the union and intersection between two uncertain sets $S_1$ and $S_2$ have the following corresponding configuration spaces:
		\begin{align*}
			\mathrm{Conf}(S_1 \cup S_2) &= \{ S_1^* \cup S_2^* : S_1^* \in \mathrm{Conf}(S_1), \, S_2^* \in \mathrm{Conf}(S_2) \}, \\
			\mathrm{Conf}(S_1 \cap S_2) &= \{ S_1^* \cap S_2^* : S_1^* \in \mathrm{Conf}(S_1), \, S_2^* \in \mathrm{Conf}(S_2) \}.
		\end{align*}
	\end{example}
	
	\begin{example}[Numerical Illustration and Uncertain Cardinality]
		Consider the universal set $U = \mathbb{R}$. Let $x_1$ and $x_2$ be two uncertain elements whose possibility domains are respectively:
		\[
		\mathrm{Val}(x_1) = \{1, 2\}, \quad \mathrm{Val}(x_2) = \{1, 3\}
		\]
		Consider the uncertain set $S = \{x_1, x_2\}$. 
		
		Assuming the states of $x_1$ and $x_2$ vary independently, the configuration space $\mathrm{Conf}(S)$ consists of $2 \times 2 = 4$ realized states:
		\begin{itemize}
			\item When $x_1 = 1, x_2 = 1 \implies S_1^* = \{1, 1\} = \{1\}$ (due to element deduplication in classical set theory), implying $|S_1^*| = 1$.
			\item When $x_1 = 1, x_2 = 3 \implies S_2^* = \{1, 3\}$, implying $|S_2^*| = 2$.
			\item When $x_1 = 2, x_2 = 1 \implies S_3^* = \{2, 1\}$, implying $|S_3^*| = 2$.
			\item When $x_1 = 2, x_2 = 3 \implies S_4^* = \{2, 3\}$, implying $|S_4^*| = 2$.
		\end{itemize}
		
		Therefore, the set of all possible configurations of $S$ is:
		\[
		\mathrm{Conf}(S) = \big\{ \{1\}, \{1, 3\}, \{1, 2\}, \{2, 3\} \big\}
		\]
		Hence, the uncertain cardinality $|S|$ has a possibility domain that constitutes an uncertain number on $\mathbb{N}$:
		\[
		\mathrm{Val}(|S|) = \{ |S^*| : S^* \in \mathrm{Conf}(S) \} = \{1, 2\} \subseteq \mathbb{N}
		\]
	\end{example}
\end{document}